\chapter{Roadmap and Open Gaps}
\label{ch:roadmap}

This closing chapter gathers the open threads left throughout both volumes into one place,
organized the same way CNA's own README states its roadmap: incremental API coverage first,
then the specific named gaps, then cross-platform strengthening.

\section{Continued API coverage, class by class}

The largest item on CNA's own roadmap is also its least dramatic: continuing to expand XNA
API coverage and behavior parity incrementally, class by class, rather than in one large
push. Chapter~\ref{ch:what-is-cna} already gave the headline numbers --- 227 of 245 public
FNA types present (92.7\%), with \texttt{.Content} (four of twelve types, by design, per
Chapter~\ref{ch:content-and-assets}) and \texttt{.Media}'s fourteen previously-shell types
(now largely closed, per Chapter~\ref{ch:media-system}'s own correction to that figure) as the
two namespaces furthest from full presence. This kind of gap closes the way most of this
book's own bug histories closed: one class, one audit, one fix at a time, tracked through the
\texttt{plan\_*.md}/\texttt{AUDIT.md} mechanism from Chapter~\ref{ch:project-practice}.

\section{The single biggest named gap: compiled .fx bytecode}

Chapter~\ref{ch:shader-fx-gap} in Volume~I covered this in full, and it remains, in the
project's own words, the single biggest real gap in CNA: no support for loading a compiled
\texttt{.fx} shader bytecode blob the way a real XNA content pipeline would produce one, with
a real, measured cost of 23 of the 86 official XNA samples in \texttt{cna-samples}
(Chapter~\ref{ch:samples-and-examples}) failing to run as a direct result. The
\texttt{ShaderEffect} runtime-HLSL-compilation path on the three Direct3D backends
(Chapter~\ref{ch:direct3d-backends}) is a real, working alternative for new or adapted
shaders, but is not bytecode compatibility, and closing this gap for real would mean building
something comparable in scope to a second MojoShader-style bytecode cross-compiler.

\section{The .xnb content pipeline: a design choice under active reconsideration}

Chapter~\ref{ch:content-and-assets} covered CNA's raw-asset-plus-JSON-descriptor content
model as a deliberate design choice, not an oversight --- and the roadmap frames a real,
binary-compatible \texttt{.xnb} reader as something the project is actively considering
rather than something ruled out permanently. Chapter~\ref{ch:models-meshes}'s account of the
real, already-shipped \texttt{.xnb} \cnaclass{Model} reader is a concrete sign of where this
is already heading incrementally, class by class, rather than as one large content-pipeline
rewrite.

\section{Named architecture-decision gaps}

Four specific gaps recur across Volume~I's backend chapters, each requiring a project-owner
decision rather than more engineering effort on its own:

\begin{itemize}
  \item \texttt{SDL\_RENDERER}'s \cnaclass{TextureAddressMode::Wrap}/\texttt{Mirror} via
        \cnaclass{SpriteBatch} (Chapter~\ref{ch:sdl-renderer}) --- blocked on choosing between
        three concrete implementation strategies, each with a real tradeoff.
  \item \texttt{SDL\_RENDERER}'s \cnaclass{Texture3D}/\cnaclass{TextureCube} construction
        (Chapter~\ref{ch:sdl-renderer}) --- blocked because the safer fix (throwing at
        construction rather than silently succeeding) has a confirmed 94-test blast radius.
  \item \texttt{EASYGL}'s real, non-\texttt{Color} \cnaclass{SurfaceFormat} GPU forwarding
        (Chapter~\ref{ch:easygl-backend}) --- blocked because it conflicts with an
        already-shipped \texttt{Texture::ValidateFormat} contract the project has already
        committed tests against.
  \item \cnaclass{Texture3D}/\cnaclass{TextureCube}'s sampler-bind architecture
        (Chapter~\ref{ch:textures-rendertargets}) --- the structural deviation from FNA
        where neither type participates in the generic texture-slot binding mechanism at
        all, so only stock effects with a dedicated field (\cnaclass{EnvironmentMapEffect}'s
        cube map) can use one.
\end{itemize}

\section{Strengthening cross-platform execution and validation}

The roadmap's final, broadest item is exactly what this Part's cross-platform chapters
covered in depth: deepening validation on Windows (moving from Wine-verified to genuinely
Windows-hardware-verified, an open item on every one of the three Direct3D backends), the
Web (moving \texttt{CANVAS} from code-reviewed to actually pixel-verified in a real browser,
Chapter~\ref{ch:canvas-ascii-backends}), and Android (broadening beyond the single real
x86\_64 emulator configuration this project currently verifies against).

\section{A closing note on how to use this chapter}

Every gap named in this chapter is named because a specific, dated source said so ---
\texttt{NEXT.md}, a \texttt{plan\_*.md} file, or a backend's own documentation, per the
practice conventions from Chapter~\ref{ch:project-practice}. As with every dated claim this
book has made, the right way to use this chapter is not to trust it indefinitely, but to use
it as a map: go to the source it cites, check whether the gap is still open, and update your
own understanding from there. That is the same discipline this project asks of itself, and
the same discipline this book has tried to model throughout both volumes.
